\section{Заключение}
    В работе предлагается подход предварительного снижения размерности перед проведением процесса кластеризации \kmeans{}. Было проведено комплексное исследование влияния методов снижения размерности PCA, Spectral Embeddings, Locally Linear Embeddings, MDS, ISOMAP, UMAP на качество кластеризации алгоритмом  \kmeans{}. Основной акцент был сделан на анализ того, как различные подходы к уменьшению размерности данных сохраняют или искажают кластерную структуру, а также как это сказывается на эффективности последующей кластеризации. В результате, было обнаружено, что применение PCA и ISOMAP на синтетических и реальных данных не приводило к снижению качества при переходе в меньшую размерность, при этом на наборе данных Digits метод ISOMAP позволил повысить качество кластеризации более чем на 10\%.
    Код проекта доступен по ссылке \url{www.github.com/alexgiving/LKMeans}.

    \paragraph{Применение суперкомпьютерного комплекса CHARISMA.}
        Исследование выполнено с использованием суперкомпьютерного комплекса НИУ ВШЭ CHARISMA \cite{charisma}.

    \paragraph{Описание применения генеративной модели.}
        В выпускной квалификационной работе использовалась генеративная модель DeepSeek-V3 Chat \cite{deepseekai2024deepseekv3technicalreport}, доступная по ссылке \url{www.deepseek.com}.
        Модель была применена для систематизации и первичного структурирования обзора литературы, что позволило выделить ключевые направления исследований и логично выстроить анализ научных источников. Кроме того, модель использовалась для оформления \LaTeX{} шаблона работы в соответствии с требованиями НИУ ВШЭ, включая корректное применение стилей, разметки и форматирования. Все сгенерированные фрагменты были переработаны и проверены автором вручную. Итоговый текст работы представляет собой результат анализа, интерпретации и авторского изложения материала.
